\documentclass{article}
\usepackage{/globals/de}
\usepackage{/globals/anki}
\ankiprefix{staatsrecht-bürokratie}
\ankitopic{Bürokratie}
\ankishow
\begin{document}
\section{Überregulierung} 
Es ist einfach die Kritik, dass es zu viele Gesetze, zu viel Staat oder zu viel Bürokratie gebe, zu finden. Diese kann in drei Teile aufgeteilt werden.
\begin{description}
 \item[Staat] Der Staat habe zu viele Aufgaben, eine zu hohe Staatsquote.
 \item[Regulierung] Es gebe zu viele Pflichten und Standards, Regulierungen wären zu komplex.
 \item[Bürokratische Organisation] Die Organisation der Bürokratie sei das Problem; Bearbeitungen seien zu langsam \etc.
\end{description}
\ankinote{arten}{Arten}{Staat (Aufgaben, Staatsquote), Regulierung (Pflichten, Komplexität), Bürokratische Organisation (Langsame Bearbeitung)}
Dass es so viele Gesetze gibt liegt unter anderem daran, dass diese in eine so komplexen Gesellschaft erforderlich sind und dass es Gesetze mehrere Ebenen gibt (\zb Gesetze auch aus dem EU-Recht).
\ankinote{gründe}{Gründe}{Komplexität, Wohlstand, EU}

\end{document}